% CA - Project Report — IEEE 2-column, ~5 pages (condensed)
% Student: Sai Shreyas Gubbi Harish | ID: 24194956 | Scalable Cloud Applications

\documentclass[conference]{IEEEtran}
\IEEEoverridecommandlockouts
\usepackage{cite}
\usepackage{amsmath,amssymb,amsfonts}
\usepackage{graphicx}
\usepackage{textcomp}
\usepackage{url}
\usepackage{hyperref}
\usepackage{booktabs}
\hypersetup{colorlinks=true, linkcolor=blue, citecolor=blue, urlcolor=blue}

\begin{document}

\title{Cloud-Native Image-to-Text Platform:\\Architecture, Implementation, and Deployment on AWS}

\author{\IEEEauthorblockN{Sai Shreyas Gubbi Harish}
\IEEEauthorblockA{Student ID: 24194956 \\
National College of Ireland --- Scalable Cloud Applications \\
Email: [your.email@student.ncirl.ie]}
}

\maketitle

\begin{abstract}
This report describes a cloud-native image-to-text (OCR) platform on Amazon Web Services. A custom REST API runs as Node.js Express on AWS Lambda behind Amazon API Gateway, with Amazon S3, DynamoDB, and SQS for storage, jobs, and asynchronous processing; OCR uses Amazon Textract with a Tesseract.js fallback, followed by post-validation (quality, blocked terms, PII-style checks). A Next.js web application (static export) is hosted on AWS Amplify (or alternatively served via Lambda/API Gateway), uses Amazon Cognito for authentication, and consumes the author's OCR API, a classmate's translation API (HTTP API), and public services including Unsplash and an AI text API. Scalability uses serverless compute, a queue-based worker, and caching. CI uses GitHub Actions with SonarCloud and OWASP ZAP; deployment is scripted with the AWS SDK. The report summarises requirements, architecture, implementation, and concise lessons learned.
\end{abstract}

\begin{IEEEkeywords}
Serverless, AWS Lambda, OCR, API Gateway, DynamoDB, SQS, Next.js, Cognito, CI/CD.
\end{IEEEkeywords}

\section{Introduction and Requirements}

\IEEEPARstart{T}{he} project delivers a full-stack cloud application: a scalable OCR \textbf{backend} (REST, JWT-auth, CRUD jobs) and a \textbf{front-end} dashboard (upload, translate, vault, optional image search and AI tools). Assignment constraints include three--five web services (author API, classmate API, public APIs), scalable backend design (FaaS, queues), and public cloud deployment with the API shared for classmates.

\textbf{Functional:} The OCR API accepts multipart or base64 JSON via API Gateway, returns text and confidence, persists jobs (UUID) in DynamoDB with per-user listing (GSI on \texttt{userId}), stores images in S3, and enforces content policies (sensitive/prohibited content must not be returned as normal text). The web app gates user features behind Cognito. \textbf{Non-functional:} scalability (Lambda, SQS, DynamoDB), JWT verification and CORS, resilience on the async path, and maintainability (separate repos/scripts).

Table~\ref{tab:services} lists integrated services. User isolation uses Cognito \texttt{sub} as \texttt{userId}; optional ElastiCache (Redis) may back the consumer for cache throughput.

\begin{table}[t]
\caption{Web services composed by the front-end.}
\label{tab:services}
\centering
\footnotesize
\begin{tabular}{@{}c l l l@{}}
\toprule
\# & Service & Type & Role \\
\midrule
1 & OCR API & Author & Upload, OCR, jobs, S3 links \\
2 & Text-to-Multiple-Languages & Classmate & Multi-target translation \\
3 & Unsplash & Public & Image search/browse \\
4 & AI text API & Public & Captions/tags (AI Tools) \\
\bottomrule
\end{tabular}
\end{table}

\section{Architecture and Design}

\subsection{Layers and flows}
The \textbf{web tier} is Next.js (App Router) with static export (\texttt{out/}). Production uses \textbf{AWS Amplify Hosting} (build via \texttt{amplify.yml} or manual \texttt{deploy-webapp-amplify.js}) with env vars for Cognito and API bases; an alternative is Lambda (\texttt{webapp-serve}) + API Gateway serving files from S3. The \textbf{API tier} is one Lambda running \texttt{server.js} (Express) behind a separate HTTP API: CORS and JWT checks in Lambda (\texttt{aws-jwt-verify}). Data stores: S3 (images), DynamoDB (\texttt{OCRJobs}, \texttt{UserS3Links}, \texttt{OCRCache}), SQS (async jobs). A \textbf{consumer} Lambda drains SQS, runs OCR, updates DynamoDB.

\textbf{Sync path:} \texttt{POST /ocr/base64} $\rightarrow$ cache check $\rightarrow$ OCR $\rightarrow$ post-validation $\rightarrow$ persist $\rightarrow$ JSON. \textbf{Async path:} \texttt{POST /ocr/base64?async=1} $\rightarrow$ pending job + S3 + SQS $\rightarrow$ consumer $\rightarrow$ complete; client polls \texttt{GET /ocr/:jobId}. \textbf{Observability:} CloudWatch logs/metrics for API Gateway, Lambdas, and SQS.

\subsection{Patterns}
\textbf{FaaS} avoids fixed servers \cite{aws-lambda}. \textbf{SQS + consumer} decouples API latency from OCR work \cite{eip}. \textbf{Cache-aside:} DynamoDB hash cache, browser \texttt{sessionStorage}, optional Redis. \textbf{Cognito JWTs} with per-user GSI queries \cite{cognito}. These meet scalability requirements better than a single EC2 monolith.

\subsection{Data model and API (summary)}
DynamoDB: \texttt{OCRJobs} (PK \texttt{jobId}; GSI \texttt{ByUserId}), \texttt{UserS3Links} (\texttt{userId}, \texttt{jobId}), \texttt{OCRCache} (content hash, TTL). Representative endpoints: \texttt{POST /ocr/base64}, \texttt{POST /ocr/base64?async=1}, \texttt{GET /ocr/:jobId}, \texttt{GET /ocr?list=mine}, \texttt{GET /users/me/s3-links}, \texttt{PUT}/\texttt{DELETE /ocr/:jobId}; \texttt{Authorization: Bearer} for protected routes.

\subsection{OCR post-validation}
After OCR, \texttt{ocr-postprocess.js} validates upload metadata (type, size), normalises text, computes \textbf{quality} metrics and score (confidence, length, noise), runs \textbf{script} analysis, applies a \textbf{blocked-word} list (\texttt{blocked-words.txt} / PII reference), and \textbf{sensitivity} checks (email/phone/SSN/card patterns and \texttt{sensitive-terms.txt}). Blocked or sensitive content yields an error response without displaying raw extracted text where policy requires.

\section{Implementation Highlights}

\textbf{Backend:} AWS SDK v3 for S3/DynamoDB/SQS. \texttt{ocr-process.js}: Textract \texttt{DetectDocumentText}, else Tesseract + \texttt{sharp}. Consumer (\texttt{api/consumer/index.js}) reuses OCR/post-process modules; deployment packages both Lambdas via \texttt{api/deploy.js} (IAM, tables, queue, API, event source mapping).

\textbf{Front-end:} Amplify Auth in \texttt{AuthProvider}; pages for upload (sync OCR), translate (classmate API), images (Unsplash), vault and settings (local storage / authenticated APIs). Client caches (\texttt{sessionStorage}, \texttt{localStorage}) reduce repeat calls.

\section{CI/CD and Deployment}

\textbf{GitHub Actions} (\texttt{webapp-ci.yml}): checkout, Node 22, \texttt{npm ci}, \texttt{npm run build}; \textbf{SonarCloud} with \texttt{npm run test:coverage} and LCOV; \textbf{OWASP ZAP} baseline against a Docker-served static build; reports in the workflow summary. \textbf{CD:} \texttt{api/deploy.js} updates Lambdas, API Gateway, and dependencies; webapp deploy via Amplify script or Amplify Console from Git. \textbf{Quality gates:} static analysis + DAST complement JWT and CORS hardening in Lambda.

\section{Conclusions}

The solution meets the brief: multi-service front-end, scalable serverless backend (SQS, DynamoDB, Lambda), Cognito isolation, and automated CI (SonarCloud, ZAP). Main lessons: (1) perceived performance depends heavily on client batching and caching, not only on hosting; (2) queues and optional caches matter for CPU-heavy OCR; (3) handling CORS and credentials in the Lambda avoids API Gateway CORS conflicts; (4) post-validation is necessary for policy and safety. Future work: dependency scanning, stricter rate limits, and broader integration tests.

\section*{References}

\begin{thebibliography}{00}
\bibitem{aws-lambda} Amazon Web Services, ``AWS Lambda,'' documentation, 2025.
\bibitem{eip} G. Hohpe and B. Woolf, \textit{Enterprise Integration Patterns}. Addison-Wesley, 2003.
\bibitem{cognito} Amazon Web Services, ``Amazon Cognito Developer Guide,'' 2025. \url{https://docs.aws.amazon.com/cognito/}
\bibitem{sonar} SonarSource, ``SonarCloud,'' 2025. \url{https://sonarcloud.io}
\bibitem{nextjs} Vercel, ``Next.js Documentation,'' 2025. \url{https://nextjs.org/docs}
\end{thebibliography}

\end{document}
